\section*{2026 年 3 月 9 日作业}
\phantomsection
\addcontentsline{toc}{section}{2026 年 3 月 9 日作业}
\noindent\textbf{截止时间：2026 年 3 月 16 日 13:30}
\setcounter{problem}{0}

\begin{problem}
已知
\[
A=
\begin{bmatrix}
2 & -3\\
3 & 2
\end{bmatrix}.
\]
求 $\cos A$ 和 $\sin A$.
\begin{solution}
容易知道
$
A=\begin{bmatrix}2&-3\\3&2\end{bmatrix}
$
是复数 $z=2+3\ii$ 的矩阵表示。定义
$
J=\begin{bmatrix}0&-1\\1&0\end{bmatrix},
$
则
$
A=2I+3J,
$
它与 $z=2+3\ii$ 是一个代数同构。

先计算
\[
\cos(2+3\ii)=\cos2\cosh3-\ii\sin2\sinh3.
\]
对应到矩阵，
\[
\cos A=\cos(2I+3J)=\cos2\cosh3\,I-\sin2\sinh3\,J
=\begin{bmatrix}
\cos2\cosh3 & \sin2\sinh3\\
-\sin2\sinh3 & \cos2\cosh3
\end{bmatrix}.
\]

同理，
\[
\sin A
=\begin{bmatrix}
\sin2\cosh3 & -\cos2\sinh3\\
\cos2\sinh3 & \sin2\cosh3
\end{bmatrix}.
\]
\end{solution}
\end{problem}

\begin{problem}
记 $z=x+\ii y$，$x,y\in\R$。关于复变量 $z$ 的函数 $f$ 可以视为关于独立实变量 $x,y$ 的二元函数，而 $x,y$ 与 $z,\bar z$ 可以互相线性表示，因此函数 $f$ 也可从形式上视为关于独立变量 $z,\bar z$ 的二元函数。证明：在这一意义下 Cauchy-Riemann 方程可表示为
\(
\partial f / \partial \bar{z} =0.
\)
\begin{solution}
设
\[
f(z)=u(x,y)+\ii v(x,y).
\]
将 $f$ 视为 $(z,\bar z)$ 的二元函数，则
\[
\frac{\partial f}{\partial\bar z}
=\frac{\partial f}{\partial x}\frac{\partial x}{\partial\bar z}
+\frac{\partial f}{\partial y}\frac{\partial y}{\partial\bar z}.
\]

由
\[
\begin{cases}
z=x+\ii y,\\
\bar z=x-\ii y,
\end{cases}
\quad\Longrightarrow\quad
\begin{cases}
x=\dfrac{z+\bar z}{2},\\[2mm]
y=\dfrac{z-\bar z}{2\ii},
\end{cases}
\]
得
\[
\frac{\partial x}{\partial\bar z}=\frac12,\qquad
\frac{\partial y}{\partial\bar z}=-\frac1{2\ii}=\frac{\ii}{2}.
\]

故
\begin{align*}
\frac{\partial f}{\partial\bar z}
&=\frac12\frac{\partial f}{\partial x}
+\frac{\ii}{2}\frac{\partial f}{\partial y}\\
&=\frac12(u_x+\ii v_x)+\frac{\ii}{2}(u_y+\ii v_y)\\
&=\frac12(u_x-v_y)+\frac{\ii}{2}(v_x+u_y).
\end{align*}
在 Cauchy-Riemann 方程
\[
u_x=v_y,\qquad u_y=-v_x
\]
下，上式即为 $0$，故 $\partial f/\partial\bar z=0$。
\end{solution}
\end{problem}

\begin{problem}
设 $m,n$ 为正整数，给定 $A\in\R^{m\times n}$。定义 $\R^n\to\R^m$ 的映射
\[
f(x)=Axx^\top x.
\]
求 $\partial f/\partial x$.
\begin{solution}
对 $x$ 做微分：
\begin{align*}
\dd f
&=\dd(x^\top x\,Ax)\\
&=\dd(x^\top x)\,Ax+x^\top x\,\dd(Ax)\\
&=2x^\top\dd x\,Ax+x^\top x\,A\,\dd x\\
&=2Axx^\top\dd x+x^\top x\,A\,\dd x\\
&=\bigl(2Axx^\top+x^\top x\,A\bigr)\dd x.
\end{align*}
故
\[
\frac{\partial f}{\partial x}=2Axx^\top+x^\top x\,A.
\]
\end{solution}
\end{problem}

\begin{problem}
设 $m,n$ 为正整数且 $m\ge n$，$A\in\R^{m\times n}$ 是给定的列满秩矩阵，$b\in\R^m$ 是给定的向量。证明：$\R^n$ 上的函数
\(
f(x)=(b-Ax)^\top(b-Ax)
\)
的最小值在
\(
x=(A^\top A)^{-1}A^\top b
\)
处取到.
\begin{solution}
\textbf{证法一：}
\begin{align*}
    f(x)&=(b^\top-x^\top A^\top)(b-Ax) \\
&=b^\top b-x^\top A^\top b-b^\top Ax+x^\top A^\top Ax \\
&=b^\top b-2b^\top Ax+x^\top A^\top Ax.
\end{align*}

这里因为 $x^\top A^\top b$ 是标量，所以可以取转置 $b^\top Ax$. 

因此
\[
\frac{\partial f}{\partial x}
=-2A^\top b+2A^\top Ax,
\]
所以在
\[
x^\ast=(A^\top A)^{-1}A^\top b
\quad(\text{因为 $A$ 列满秩})
\]
处有驻点。又
\[
\frac{\partial^2 f}{\partial x^2}=2A^\top A\succ0
\quad(\text{同样因为$A$ 列满秩}),
\]
故 $f(x)$ 严格凸，最小值存在且唯一，于是 $x^\ast$ 即唯一最小值点。

\textbf{证法二：}将 $b$ 分解为落在 $A$ 的像空间上的部分 $b_1$ 与垂直部分 $b_2$：
\[
b=b_1+b_2,\qquad b_1\in\mathcal{R}(A),\ b_2\in\mathcal{R}(A)^\perp.
\]

从直观上而言，
\[
(b-Ax)^\top(b-Ax)=\|b-Ax\|_2^2
\]
就是 $b$ 到 $\mathcal{R}(A)$ 的距离平方，其最小值在 $b$ 在 $\mathcal{R}(A)$ 上的投影点处取到.

此时有 $b_1 = Ax$,\quad $b_2 = b - Ax \perp \mathcal{R}(A)$.

即：

\[
A^\top (b - Ax) = 0 \Leftrightarrow x=(A^\top A)^{-1}A^\top b.
\]

\end{solution}
\end{problem}

\begin{problem}
构造非奇异上三角矩阵 $P$ 使得
\[
P^{-1}
\begin{bmatrix}
1 & 1 & 5 & 6\\
2 & 1 & 7 & 8\\
0 & 0 & 2 & 3\\
0 & 0 & 0 & 4
\end{bmatrix}
P=
\begin{bmatrix}
1 & 1 & 0 & 0\\
2 & 1 & 0 & 0\\
0 & 0 & 2 & 0\\
0 & 0 & 0 & 4
\end{bmatrix}.
\]
\begin{solution}
设
\[
P=
\left[
\begin{array}{c|c}
I&A\\ \hline
0&B
\end{array}
\right],
\]
其中 $A,B\in\R^{2\times2}$。题设等价于
\[
\left[
\begin{array}{cc|cc}
1&1&5&6\\
2&1&7&8\\ \hline
0&0&2&3\\
0&0&0&4
\end{array}
\right]
\left[
\begin{array}{c|c}
I&A\\ \hline
0&B
\end{array}
\right]
=
\left[
\begin{array}{c|c}
I&A\\ \hline
0&B
\end{array}
\right]
\left[
\begin{array}{cc|cc}
1&1&0&0\\
2&1&0&0\\ \hline
0&0&2&0\\
0&0&0&4
\end{array}
\right].
\]
按分块乘法得到
\begin{align}
\begin{bmatrix}1&1\\2&1\end{bmatrix}A+\begin{bmatrix}5&6\\7&8\end{bmatrix}B
&=A\begin{bmatrix}2&0\\0&4\end{bmatrix},\tag{1}\\
\begin{bmatrix}2&3\\0&4\end{bmatrix}B
&=B\begin{bmatrix}2&0\\0&4\end{bmatrix}. \tag{2}
\end{align}

对 (2)，有
\[
\operatorname{spec}\!\begin{bmatrix}2&3\\0&4\end{bmatrix}
=\operatorname{spec}\!\begin{bmatrix}2&0\\0&4\end{bmatrix},
\]
所以有无穷多组解，取其中一组
\[
B=\begin{bmatrix}1&3\\0&2\end{bmatrix}.
\]

对 (1)，由于
\[
\operatorname{spec}\!\begin{bmatrix}1&1\\2&1\end{bmatrix}
\cap
\operatorname{spec}\!\begin{bmatrix}2&0\\0&4\end{bmatrix}
=\varnothing,
\]
故解唯一。解得
\[
A=
\begin{bmatrix}
-12&\dfrac{118}{7}\\
-17&\dfrac{165}{7}
\end{bmatrix}.
\]

因此可取
\[
P=
\begin{bmatrix}
1&0&-12&\dfrac{118}{7}\\
0&1&-17&\dfrac{165}{7}\\
0&0&1&3\\
0&0&0&2
\end{bmatrix},
\]
它是非奇异上三角矩阵，并满足题设关系。
\end{solution}
\end{problem}

\begin{problem}[选做]
设 $n$ 为给定的正整数，并记
\(
\mathcal{D}_n=\{X\in\R^{n\times n}:\det(X)>0\}.
\)
在 $\mathcal{D}_n$ 上定义函数
\(
f(X)=\ln\det(X),
\)
求 $\partial f/\partial X$.
\begin{solution}
这是一个复合函数。令
\[
z=\ln u,\qquad u=\det(X).
\]
则
\[
\dd z=\dd(\ln u)=\frac1u\,\dd u=\frac1u\,\dd(\det X).
\]

又因为行列式可按行展开， $\dd (\det (X)) = x_{i1}A_{i1} + x_i2A_{i2} + \ldots + x_{in}A_{in}$,

因此
\(
\dd(\det (X))
=
\begin{bmatrix}
A_{11}&A_{21}&\cdots&A_{n1}\\
A_{12}&A_{22}&\cdots&A_{n2}\\
\vdots&\vdots&\ddots&\vdots\\
A_{1n}&A_{2n}&\cdots&A_{nn}
\end{bmatrix} \dd(X)
=\operatorname{adj}(X) \dd(X)
=\det(X)X^{-1} \dd(X).
\)

于是
\[
\dd z
=\frac1{\det(X)}\,\det(X)X^{-1}\dd X
=X^{-1}\dd X.
\]

因此
\[
\frac{\partial f}{\partial X}=X^{-\top}.
\]
\end{solution}
\end{problem}

\begin{problem}[选做]
考察函数
\[
f(\lambda)=
\begin{cases}
\exp(-\lambda^{-4}), & \lambda\ne0,\\
0, & \lambda=0.
\end{cases}
\]
证明：$f$ 不是解析函数，但在复平面上处处满足 Cauchy-Riemann 方程.
\end{problem}

\begin{problem}[选做]
构造次数最低的一元有理系数非零多项式 $p(\cdot)$，使得
\[
p(\sqrt2+\sqrt3+\sqrt6)=0.
\]
\end{problem}
